\documentclass[10pt,a4paper]{amsart}
\usepackage[margin=19mm]{geometry}
\usepackage{amsmath,amssymb,mathtools}
\usepackage[scr=rsfs,cal=euler]{mathalfa}
\usepackage{xfrac}
\usepackage[unicode,hidelinks,bookmarks=false]{hyperref}
\newcommand{\ie}{i.e.\ }
\newcommand{\cf}{cf.\ }
\newcommand{\resp}{respectively\ }
\newcommand{\Subsection}{\subsection}
\providecommand{\nobreakdash}{\nobreak\hbox{-}}
\setlength{\emergencystretch}{2em}
\multlinegap0pt
\usepackage{bm}
\usepackage{bbm}
\usepackage{dsfont}
\usepackage{xspace}
\usepackage{cancel}
\usepackage{mathtools}
\usepackage{amsthm}
\newtheorem{theorem}{Theorem}
\newcommand{\R}{\mathbb{R}}
\title[Multiplicity bounds for nonzero upper-Laplacian eigenvalues ]{Multiplicity bounds for nonzero upper-Laplacian eigenvalues of simplicial complexes}
\author{Vinayak Gupta}
\date{Source: arXiv:2609.29677v1 (CC BY 4.0)}
\begin{document}
\maketitle

\noindent\emph{Source and licence.} Multiplicity bounds for nonzero upper-Laplacian eigenvalues of simplicial complexes. Version arXiv:2609.29677v1 (CC BY 4.0), distributed under the Creative Commons Attribution 4.0 International licence, as recorded in the arXiv licence metadata for this exact version and shown on the abstract page https://arxiv.org/abs/2609.29677v1. Full rights notice: licenses/arxiv-2609.29677-CC-BY-4.0.txt.\par\smallskip

\noindent\textit{Editorial scope.} Counted unit: the Theorem whose source label is \texttt{thm:simplicial-faria} in the article (source0.tex lines 1023--1028, source label \texttt{thm:simplicial-faria}), delivered together with the article's own complete original proof of it (source0.tex lines 1030--1100, one \texttt{proof} environment of 71 lines). No mathematics is written, repaired or re-derived: statement and proof are the article's own text line for line. Required context retained uncounted: none. Every definition, display and label of those lines is retained unchanged. Omitted from this package and not redistributed: 1--1022, 1029--1029, 1101--1234. Nothing inside a retained excerpt is omitted or altered: every retained body is delivered exactly as the article's lines are written, so no omission receipt of any kind accompanies this package. Citations: the retained excerpts carry no citation command at all, so no bibliography entry of the article is transcribed and no reference is redistributed. Reference adaptation: none was needed and none was made; every reference inside the delivered unit resolves inside it, so no unresolved reference or citation is produced. Printed slips kept literal: no printed word, symbol, hypothesis or number of the counted statement or of its proof was altered. Layout and class: the article's own class is \texttt{article} with packages including fontenc, geometry, amsmath, amssymb, amsthm, mathtools, xcolor, microtype, float, tikz, hyperref; the wrapper is the lane's pinned \texttt{amsart} editorial wrapper at 10pt on A4, which restates the theorem-like environments the retained text needs and adds the article's own macro definitions that the retained text needs (\texttt{\textbackslash R}), with extra wrapper packages (bm, bbm, dsfont, xspace, cancel, mathtools). The delivered numbering is the unit's own. Assets: no retained span contains a figure environment, a graphics-inclusion command, a PDF-page inclusion, a table or a TikZ picture; no image file is copied into this package, the package declares no asset dependency and no image file is redistributed with it. Licence: version arXiv:2609.29677v1 of this article is distributed under the Creative Commons Attribution 4.0 International licence, as recorded in the arXiv licence metadata for this exact version and shown on the abstract page https://arxiv.org/abs/2609.29677v1. A full rights notice is delivered at licenses/arxiv-2609.29677-CC-BY-4.0.txt. Characters: every retained excerpt is pure ASCII, so the delivered engine, XeLaTeX with its default Latin Modern text font, reports no missing character.
\begin{theorem}\label{thm:simplicial-faria}
Let $\Upsilon$ be a finite pure $d$-dimensional simplicial complex, where $d\ge1$.  Then
\[
m_{\Upsilon}(d)\ge p_d(\Upsilon)-q_d(\Upsilon).
\]
\end{theorem}

\begin{proof}
Choose orientations of the facets and ridges of $\Upsilon$, and put
\[
B=\partial_d,\qquad M=B^{\mathsf T}B.
\]
Since $d>0$, the multiplicity of $d$ as an eigenvalue of $M$ equals its multiplicity as a nonzero eigenvalue of
\[
BB^{\mathsf T}=L_{d-1}^{\mathrm{up}}(\Upsilon).
\]

For each quasi-pendant ridge $R$, let
\[
\mathcal P_R
=\{F\in\Upsilon_d:F\text{ is pendant and its unique nonfree ridge is }R\},
\qquad
t_R=|\mathcal P_R|.
\]
The sets $\mathcal P_R$ partition the pendant facets, and therefore
\[
\sum_R t_R=p_d(\Upsilon),
\]
where the sum is over the $q_d(\Upsilon)$ quasi-pendant ridges.

For $F\in\mathcal P_R$, let $\varepsilon_F\in\{1,-1\}$ be the incidence sign of $R$ in $\partial_dF$, and define
\[
U_R=
\left\{
x\in\R^{\Upsilon_d}:
\operatorname{supp}(x)\subseteq\mathcal P_R,\quad
\sum_{F\in\mathcal P_R}\varepsilon_Fx_F=0
\right\}.
\]
The displayed relation is one nonzero linear equation on the $t_R$ coordinates, so
\[
\dim U_R=t_R-1.
\]

We claim that every $x\in U_R$ satisfies $Mx=dx$.  If $F\in\mathcal P_R$, then the only ridge of $F$ that can belong to another facet is $R$.  Since a $d$-simplex has $d+1$ ridges,
\[
\begin{aligned}
(Mx)_F
&=(d+1)x_F+
\sum_{\substack{G\in\mathcal P_R\\G\ne F}}
\varepsilon_F\varepsilon_Gx_G\\
&=(d+1)x_F+
\varepsilon_F\left(-\varepsilon_Fx_F\right)
=dx_F.
\end{aligned}
\]

Now let $H\notin\mathcal P_R$.  If $H$ does not contain $R$, then it shares no ridge with any facet in $\mathcal P_R$, and hence $(Mx)_H=0$.  If $R\subset H$, and $\varepsilon_H$ denotes the incidence sign of $R$ in $\partial_dH$, then
\[
(Mx)_H
=\varepsilon_H\sum_{F\in\mathcal P_R}\varepsilon_Fx_F
=0.
\]
Since $x_H=0$, both cases give $(Mx)_H=dx_H$.  Thus
\[
U_R\subseteq\ker(M-dI).
\]

For distinct quasi-pendant ridges $R$ and $S$, the sets $\mathcal P_R$ and $\mathcal P_S$ are disjoint.  Hence the spaces $U_R$ have disjoint supports and their sum is direct.  Consequently,
\[
\begin{aligned}
m_{\Upsilon}(d)
&\ge\sum_R\dim U_R\\
&=\sum_R(t_R-1)\\
&=p_d(\Upsilon)-q_d(\Upsilon).
\end{aligned}
\]
\end{proof}

\end{document}
